\documentclass[12pt]{article}
\usepackage[T1]{fontenc}
\usepackage[utf8]{inputenc}
\usepackage[brazil]{babel}
\usepackage[margin=1in]{geometry}
\setlength{\parindent}{0pt}
\begin{document}
\section*{Tema 44}
Processo: PEDILEF 0021992-38.2008.4.01.3600/ MT\\
Situacao: Julgado\\
Ramo: DIREITO ADMINISTRATIVO\\
Questao: Saber qual a natureza jurídica da Gratificação de Estímulo à Docência no Magistério Superior (GED) após a MP n. 208/2004 - se parcela remuneratória geral ou "pro labore faciendo" - para fins de equiparação da pontuação dos servidores ativos aos inativos.\\
Tese: "A Lei 11.087/05, resultante da conversão da Medida Provisória 208/2004, não modificou a natureza pro labore faciendo da GED, porquanto trouxe apenas alteração nos pontos a serem atribuídos a ativos e inativos, preservando-se a diferenciação estabelecida na Lei 9.784/1998, inclusive quanto aos servidores docentes cedidos." (PET 9600/STJ) Obs: O referido entendimento foi aplicado também à PET 9645/STJ, oriunda deste representativo.\\
Fonte: https://www2.jf.jus.br/phpdoc/virtus/mostradocumento.php?Process=00219923820084013600\&seq=3
\end{document}
